My forbears, the old men and women folk, were of the Mexican People Clan. With great hardship they came back from the place called Fort Sumner, packing their goods on their backs, and afterward lived here at Lukachukai. Here they began life anew. Then from their home at a place called Spring in the Clay they started moving over and across yonder mountain. While they were making this move they stopped for the night, and it is said that I was born where they halted.

Let's see, I believe it was in the month of October I was born on a night when a heavy snow had begun to fall. My grandfather, a man by the name of White Mexican Clansman, and another man called The Son of Former Sweathouse, were the only men. They and several womenfolk worked hard on me.

My grandfather had a fine wether goat that he was quite proud of. He killed it there, and the people feasted with great joy. Then on the day following the night of my birth we again started to move. In those days people went about only on horseback or afoot.

On the day following the night on which she gave birth to me, my mother went on foot as far as the people moved that day. And since it had snowed the night before they had to go through the snow. It is remarkable how risky life was in those days. Nowadays when a woman gives birth they take special care of her, telling her not to catch cold. Back at that time it was different.

As time went on I somehow grew up. I do not remember much about it. I do not remember in detail much of what happened before I was about ten years old. When morning would come and I would be sleeping nicely, my grandfather would throw me out of bed.

``Come on! Come on! Why are you lying down? Go run a race. If you're weak the first harmful thing that comes along will run over you. If you are strong you will lie hudled in death only after the dirt is torn and furrowed around you (i.e. you will give up only after a struggle),'' my grandfather used to tell me. The old folks all used to say that.

They probably said this because they knew of the wars that were going on. That is the way my grandfather drove me about and spoke to me. He did not speak in anger. He did not speak harshly. He told it to me in a nice way. So that is the way I spent my days from that time on.

After I had reached the age of eighteen I did various kinds of odd jobs. I had an uncle by the name of Forked House. He was a young man, and he adopted me. He started sending me on chores. I would take care of his horses for him. He had a wagon and a team. He would tell me to haul wood with it, or to work with it.

At the time I was about eighteen years old the person who started the present trading post at Blue Clay Point moved there -- I used to know just when it was. They used to call the Former Old Interpreter (Chee Dodge) Chee (Red) when I first heard of him. Chee was a good friend of a white man from Ugly House (Manuelito, N.M.) known as Big Lump Setting Up (S.E. Aldrich). Finally the two of them built the trading post at the place called Blue Clay Point. It was in the Spring when they moved there. There was an epidemic in progress, and we now refer to it as ``the time when the throat killed many.'' It was at that time, in the spring, that they moved to Blue Clay Point, it was said. That is how it came about that the one who would be called The Interpreter (Chee Dodge) and the one who would be named Big Lump Setting Up came to move there.

They built the trading post and four years later perhaps, I do not recall exactly when, the one I called my uncle told us to haul wood there. They (the traders) probably called for the wood.
There was a boy of the Big Water Clan with whom I always chummed around. We started doing that kind of work. We went about calling each other ``younger brother,'' and we hauled wood to the trading post. We heard reports that there were people going about to get children put into school. It was said that they were then over at Chinle, and we heard that they were coming here sometime soon. We would discuss it when we took the wagon after wood.

``Say, let's go to school,'' I said to my companion.

``All right,'' he said.

And we went about, talking enthusiastically about it. In the evening after we had unloaded another load of wood we started back home in the wagon. When we got back to the hogan we found that my uncle was there. And it turned out that my grandfather, who was called Big Mexican Clansman, was there too. So we told them of our plan. ``They are putting chidlren in school, it is said.'' We told them how we had both said, ``Let's go to school.''

``It's probably all right. It's fine indeed, my children. A party is in its way coming to us. You cannot escape them anyway,'' my uncle and my grandfather said to us.

On the next day we went to the trading post, and my grandfather put his thumbprint on a paper for us. It was reported that the children from Chinle who were being taken to school would get there in two days, and we heard the grownups talking to one another about it.

Of those menfolk who were talking to one another about it, one was called Weela. He lived over behind Round Rock and was one of the leaders at that time. One was called Bead Clan Gambler, and he was a member of the police force. He was a very strong, husky man. He was from Chinle, but he was a member of the police force at Fort Defiance. He had come on ahead with the white man from Fort Defiance called Little Chief (Agent Dana Shipley). It was there that my grandfather spoke and promised us.

``I am putting these two boys in school.'' he said by way of promise.

Since it was reported that it would be another day before the wagon load of children would arrive, we went back home. We spent the night there and then returned on foot.

And we found that the children had arrived. There was some man of the Bitter Water Clan whom people called Pug Nose. He was going around with the party, taking care of the children. He knew a little English, and that was probably why he was told to take care of the children. He was busily running around the children who had been hauled in, chattering in English. He was keeping them together. The children from Blue Clay Point who were to be put in school had not all been brought in as yet. It was said that the wagon with the children would start off in two days. They were to go by way of Lukachukai, Tsaile, and Crystal, and on over to Fort Defiance. That was probably their plan.

But the people who live over beyond yonder mountain in the area of Cove and Red Rock were thinking ugly thoughts about the plan in connection with education. It is just recently that these people became old men, and they have all died. At that time, when they acted as they did, they were young men. They said that if this business of taking children to school got to them they would really do something about it. As we came to find out, they were saying that (the school) would not get a one of their children. A man by the name of Black Horse was the leader of that faction. They had no doubt heard that in a certain number of days a party would come to Blue Clay Point for the purpose of putting children in school. So they began to think of the matter with bitterness.

``Boys, let's go there. Come on!'', they probably said. So they probably set out from there for this purpose (of opposing the placing of children in school). We had not heard of this and knew nothing about it. They met (the school party) at Blue Clay Point. 

Near the trading post there was a house which at that time, we called the Ugly House. The house served the purpose of providing a place to sleep for the people who came to the trading post. Black Horse's party went over there, and they no doubt came with evil thoughts. The people who lived hereabout had heard nothing of their plans but at night they spoke about themselves, revealing their intentions.

The one I referred to as my grandfather, and another uncle of mine called Little Boy, apparently went there at night and heard them telling about it. Among those present were the men known as Black Horse, Limper, Tall Red House Clansman, Old Bead Chant Singer -- who was of the Red House Clan -- and his cousin, a Bitter Water Clansman called Bead Chant's Nephew, and another man with whom I was merely acquainted, and who was a relative of mine by marriage. His name was Yellow Fermentive Chewer. These are the only ones I can remember of those who came from there (on the other side of the mountain). I do not remember who the rest were. They probably planned to go in to see the Agent on the next day. They no doubt told what they thought -- what their opinions were regarding this school business that had started in connection with us. They told these things to the people from here who had gathered. So people began to think along the lines of their planning. ``Now men, is there anyone here who can do it (a chant)? Long ago, people of old had a story of some kind of chant called `Talk One Into the Grave.' Who of you knows it?'' said Black Horse, asking that it be perfomed. 

These people from the other side of the mountain were saying this. This is what that is what I heard. I wasn't at the meeting myself. And I don't know just how this chant goes.

``I do,'' said the one I referred to as Little Boy.

``Two of you are needed for it,'' said Black Horse.

So then my grandfather, the man I spoke of as Big Mexican Clansman, volunteered to join him to carry on the ceremony. South from the trading post there is a ruin which we call in Navajo ``Shattered House.'' Someone burned it long ago. It is said that they were Anasazi. It is black there like ashes. It was there that they carried on the ceremony. I don't know how it was done. That is what took place that night.
As for us, we spent another night at my uncle's hogan. Today we had been placed among children who were to be taken off to school. On the morrow we would start off with them. So for that purpose, in the morning we again started off from the hogan afoot. We left with menfolk, with the one who would thumbprint the papers for us and who would vouch for us. It was my grandfather, Big Mexican Clansman. He stood up for the two of us.

When we got to the store we found that many people had assembled. Many horses were standing about. At that time horses were the only means people had for transportation. Some of us did not know what had been done the night before. Someone by the name of Black Horse had brought a party from the other side of the mountain. That is all that we found out.

We were told that now there would be work here making out papers for more of our children. There were three Navajo policemen there in that connection. One of them turned out to be Bead Clan Gambler, one was Singed Man from Fort Defiance -- he was also known as Son of Former Rag Man. Another was Interpreter's (i.e. Chee Dodge's) brother-in-law, a Red Streak Into Water Clansman. This man was killed recently at St. Michaels at a rodeo. He was killed by a racehorse on the track during the race. At that time he was a policeman. I can't think what his name was -- I merely knew him by sight. He was merely called Interpreter's Father-in-law: He was a policeman back at that time. So there were three policemen. It happened that way. Then we were told that the time had come. We saw the people going inside the trading post, so we just went with the crowd. We two who were going to school stayed together. As for the rest of the children who were going to school, I don't know anything about them.

At the time we had brought the last load of wood; while we were unloading it we noticed two women leading a horse nearby. They were probably on their way home. One of these women had promised her son for school. [His name was Toohnii and she had promised him for school] The woman glanced our way. Her name was Red Buttock Clan Woman. She is still living and her hair is white. As she passed by she said, ``It is said that those boys who are unloading wood over there have been promised for school -- they will go to school too. Golly, why they are big enough to be of use here now'', she said of us. I don't know what she had in mind in so saying. We just laughed about it afterward. But that was right. We had volunteered for school. We were going to get an education. We considered school as good news. That was what I thought. I always told my friend about what I myself was thinking. ``You see, it's something to strive for and covet. Look at Chee's trading post here. Quite some time ago when they moved back from Fort Sumner he had a hard time. He started his life by poking around in the horse manure in a stable. He used to be out there early in the morning with a stick especially made for that purpose. He could be seen sitting out there silhouetted against the sun, with the light shining red through his hair. This is how he came to be called Chee (Red). Then some white people took him in. That is how he got his start in life. That is how his mind began to develop. He got ahead and now here he is in his enviable position. He knows English'', I said. That is what I had an ambition to do at that time. It was for that reason that I volunteered myself for school. And the person who had joined me thought likewise. With this in mind I wanted an education badly then. But it was destined to turn out differently.

The people went into the trading post. It was packed full. Over here on one side the counter ran. Further back in the room was a swinging gate. It was out through there that Chee came. The brother of The One Who Has Eyeglasses (John Lorenzo Hubell) from Ganado, a man called the Bat, was a clerk there. He was over to one side behind the counter. The Bat was a man who could understand Navajo very well. A little later the one called Little Chief (Shipley) came out and stopped beside Chee. Chee was his interpreter as he began to make a speech to the people. Black Horse was standing against the counter over to one side. The people of his band were standing with him. The people who had promised their children for school were named, and we were told many things about the school.

Then Black Horse spoke up and said, ``This business of taking children away from people to put them in school -- when is it going to affect the people from over beyond the mountain?'', [he said].

``It will reach you sometime. Tomorrow these will start out and will be routed right along the mountainside'', said Little Chief (Shipley)

``We'll not give you a one of our children. And we'd just as soon fight over the matter as not'', said Black Horse balking stubbornly.

Speaking this way to each other the Agent and Black Horse exchanged many words.

The one I referred to as Limper was standing near Black Horse with a blanket wrapped around his middle. And those of this band were lined up, one behind the other.

``Come on, you boys. Remember what you said,'' said Black Horse.

The one called Limper was the first to hop in there, and he grabbed the Agent by the collar. Then they all rushed in. Chee jumped over the gate at the back of the room, and chaos followed.
``Outside with him!'' voices were saying.

They started out with him. As they were taking him outside I crawled and squeezed myself out among them. Just then they locked the door from the inside. Two of the boys who were going to go to school were locked in. The one I referred to as River Man, who is now known as Short Hair, and another boy with him. These were the ones who were locked inside. On account of this fact the man I spoke of as Weela took the Agent's side -- it was because one of the boys locked in there was his son.

The mob was carrying the Agent away. Not far from there, there was a drop. There was a wash in the blue clay with a point of land on either side. That is how the trading post got its name. It was a long drop.
``Throw him down there! [Down there, down there!]'' voices were heard saying.

A lot of people were standing alongside the trading post, and I among them. My uncle was standing beside me, and I don't remember who else was standing there. As the people carried the Agent along they beat him with their fists. They were beating him up. But as they carried him further away the one called Bead Clan Gambler went running from here where we were standing.\footnote{The English translation in \citet[27-29]{Young:1952cy} continues here with approximately two pages of narrative describing additional events surrounding the fight, including the most violent scenes. Because no corresponding Navajo text is provided in the original source, these passages have been omitted from the present edition. The translation resumes at the next point where the Navajo original is available.}

Over to one side of Black Horse the people were merely milling around. All of them had their horses standing there nearby. Many more joined him, and among these were some women. These women too had spent the night there with them. They had stayed to cook for the men. One of these women had come from over beyond the mountain. I don't know exactly where she came from.

When this happened Black Horse's henchmen spoke without effect because there were others who were trying to keep order and who said ``No'' to the proposal to attack the soldiers.

I was one of those who came back (on the next day) to where the people were. And I found that the people from beyond the mountain -- those among whom the news had been spread, and who were sent for -- had arrived. They had guns and bows and arrows. These were many people. There were also many of us who had joined in just as spectators.

A man of the Salty Water Clan whose name was Slender Lava had his family there, and his wife cooked all night for the men who were making the trouble. In the morning the women again prepared food for them. Then the people noticed that the soldiers were moving along over there.

``All right now! All right! We'll take them right over there in the flat. All right now! Before they move in! Before they move in! Let's get them!'' said Black Horse dashing about but without effect.

But his boys were really impatient. They were rubbing their guns. Then the one to whom I referred as Slender Lava said, ``No!'' He spoke thus because he knew they lacked the wherewithal to win the fight. A man of the Water's Edge Clan, called Sucker, also said, ``No!'' And of the men from hereabout with whom I came and who I mentioned before, my uncle joined in saying, ``No.''

``Wait. Wait. Speak no more like that, my Elders, my friends,'' people said, opposing them.

Then the person whom I said was known as Sucker was sent to the trading post. The soldiers moved into the trading post. They had pack mules, And there was a line of these animals far into the distance. I don't remember what the number of the soldiers who arrived was said to be. I can't recollect. But when they had moved into the past, Sucker was sent over there. Presently he came back from there on horseback.

``They say to wait, boys. They say to wait,'' she said.

Then Black Horse's boys really became impatient. They sent Sucker back over there again.

``Go over there now and tell them, `No'. Tell them not to give us that kind of talk.

Tell them it's going to be now. We'll tear the building to pieces. What can they do? Go tell them that's our answer,'' said Black Horse.

Sucker galloped over there and told them that, and then he returned at a gallop.

``They tell you to wait. The one who says that is the leader of those solders who arrived and he is a War Chief (Officer). It is not the man you threw out yesterday who is saying that,'' he said as he ran back again on horseback.

Then the men over on this side said, ``It's a fact. You never know. You don't know what they have in mind in saying this -- they might have a better idea for settling it. So don't say anything,'' said those on this side who were acting as go-betweens.

Finally many of us started over there, some on foot and some on horseback. We got to the trading post. The soldiers were there on the inside. There were many pack mules, more than twenty I guess. They had all been inside the wool storage shack, with their packs still on. Across from the wool storage shack a gun had been placed to cover its approaches. So in case Black Horse carried out his threat, the soldiers were ready.

Sucker rode over there to where Black Horse and his band were.

[`Now you can come over',] the War Chief (Army Officer) says, [...] `but come slowly, slowly','' said Sucker. 

That's what he reported to them. Then across and opposite us on the other side of the pace called Blue Clay Point horses were seen to be going, and many people too. There were more than twenty. They could be seen going up there over the crest of the hill. People watched and wondered if they had just decided to go back home. But they moved on over the hill and then turned and started coming in this direction along its base. Over here on this side the soldiers were all prepared and were inside the trading post. The entrance way was large, and on both sides of it soldiers were stationed. Over there on this side the people were impatient to shoot. They were really stroking their weapons. They were rubbing their guns, it is said. They were really foolhardy -- at least that's what the spectators were saying.

The white man called The Bat used to say that he was on the soldier's side ``I too had my gun ready like this,'' he used to say.

People would laugh at him when he told about this. He was a trader.

``If they appear over the hill -- if they're serious in their threat -- if they come close, I'll go out there in the open to meet them,'' said the War Chief (Officer). ``If they are serious about it, and if one of them aims at me with one of those guns you see them holding up out there, all of you fire a volley at him,'' said the officer. That was the plan on the inside.

Then they appeared, coming over the hilltop. They were coming toward us. At this point Sucker rode back and forth telling them to take it easy. That was the one called Sucker. Later, when his voice gave out, he merely gestured about. We were standing by the trading post. There were many of us. We were standing close to a recess in the wall of the building where we could run for protection. When they had drawn near, [...] [he] brought out a chair, [(the one called) The Old Interpreter]. He put it down beside the building, and then went back inside. We were all looking over toward where the horses were approaching. They were holding many guns up, and they were also holding arrows. There were none without weapons. There were many coming, and they were spread out quite wide.

Then the white man came out. It was the War Chief. He was from Fort Wingate and was dressed in a black coat. He took it off and draped it over the chair that had been placed there. Then he unbuttoned his cuffs. He rolled up his sleeves like a person who is going to start washing himself. When he came out he came with his arms up and his hands behind his head. Then Chee Dodge came back out. [While the people were only looking over there], The Old Interpreter came back out. ``He says to come closer,'' said Chee to the horsemen. They drew closer. ``Come closer,'' he again said. So the horsemen drew closer ``Now that's close enough, he says to you, my friends. That's enough, Black Horse. Lay your gun on the ground and go up to him, the officer says, said Chee, referring to Black Horse as some kind of relative like ``older brother.''

Several of them then got quickly down from their horses. They came up one after another. The officer still stood with his hands raised. Black Horse walked up and shook hands with him. I don't remember what he said at that time. It was Black Horse. What the officer said was interpreted. I still remember some of that. ``He says to you, `fAll right, that's fine','' said Chee to Black Horse, interpreting. ``That's enough, enough'', he said to him.

Then those who followed him came up one by one and shook hands with the officer. ``It's enough'', said Black Horse and the officer to each other, as they embraced and patted each other on the shoulder.
That's all I remember. At that point peace was restored, I guess. Those who were behind Black Horse all shook hands with the officer. They did this after laying their weapons on the ground. They came up and did it empty-handed. That's how peace was restored. And there was the sound of people expressing their thankfulness.

``That's enough my boys. Now food will be passed out for you. You can go back home braced by a meal,'' they were told.

With Chee Dodge interpreting, that's how we heard it. So peace was then restored.

I didn't realize the seriousness of it at that time. At that time I was just a ``young punk.'' that's probably the reason. I was just there with the crowd for fun. I had fun with it, just like at something that is carried on for fun.